% !PRINTED 269
\weberpage{269}
%<<CONT
Substitution dritter Art aus $\Lambda_p$, so können wir die Substitutionen
 $U$, $V$ aus $L_p$ oder aus $\Lambda_p$ so bestimmen, dass
\[
A = U\,S_2^r\,U^{-1}, \qquad A = V\,S_3^r\,V^{-1}.
\]
Setzen wir also, dem Werth $r=1$ entsprechend,
\[
A_1 = U\,S_2\,U^{-1}, \qquad A_1 = V\,S_3\,V^{-1},
\]
so ergiebt sich
\[
A = A_1^r, \qquad A = A_1^r,
\]
und $A_1$ kommt in $L_p$, $A_1$ in $\Lambda_p$ vor.

Ist ferner $B$ eine Substitution erster Art aus $L_p$, so kann
man wieder, wenn man
\[
T_1 = \begin{pmatrix} 1, & 1 \\ 0, & 1 \end{pmatrix}
\]
setzt, $U$ aus $L_p$ so bestimmen, dass
\[
B = U\,T_1^r\,U^{-1} = B_1^r
\]
wird. Hieraus ergiebt sich der Satz:

\gesperrt{Man kann in der Gruppe $L_p$ (oder $\Lambda_p$) die Substitutionen
 der ersten, zweiten und dritten Art mit Hinzuziehung
der Identität in Cyklen von $p$, $\dfrac{p-1}{2}$, $\dfrac{p+1}{2}$
Gliedern anordnen.}

Zwei solche Cyklen können ausser dem Einheitselemente
kein Glied gemein haben.

Dies ist zunächst evident bei den Substitutionen der ersten
Art, deren Grad gleich $p$ ist; denn bei diesen lässt sich der
ganze Cyklus durch Potenzirung aus einem beliebigen von der
Einheit verschiedenen Elemente ableiten.

Wenn aber in zwei Cyklen der zweiten oder dritten Art ein
gemeinsames Element vorkommt, so kann die Gruppe so transformirt
 werden, dass die beiden Cyklen die Gestalt bekommen:
\weq{12}{\begin{aligned} &1, \quad U\,S\,U^{-1}, \quad U\,S^2\,U^{-1}, \ \ldots \\ &1, \quad S, \quad S^2, \ \ldots, \end{aligned}}
worin $S = \begin{pmatrix} \sigma, & 0 \\ 0, & \sigma^{-1} \end{pmatrix}$ die Normalform hat.

Ist nun für irgend zwei von Null verschiedene Exponenten
$r,t$
\[
S^r = U\,S^t\,U^{-1},
\]
%CONT>>
